% GENERATED FILE. Edit canonical lessons, then run npm run generate:books.
% canonical-source-hash: fnv1a64:94e0d08d888d72db
% canonical-lessons: KA-C325-vivaha, KA-C325-arataksate, KA-C325-varadaksine, KA-C325-harasu, KA-C325-dibbana

\chapter{A Wedding Is Planned}
\label{ch:ka-a-wedding-is-planned}
\hlchaptermodality{325}

\begin{chapteropening}
\textbf{By the end of this chapter:} \emph{I can say a marriage, a wedding reception, a dowry, to bless and a wedding party.}

Complete the last lesson of chapter 325: I can say a marriage, a wedding reception, a dowry, to bless and a wedding party.
\end{chapteropening}

\section[\texorpdfstring{\kn{ವಿವಾಹ} (vivāha)}{ವಿವಾಹ (vivāha)}]{\kn{ವಿವಾಹ} (vivāha) — a marriage}

\label{lesson:KA-C325-vivaha}



\begin{quote}
Before the new one: say the Kannada for to extinguish, then the Kannada for to rescue.
\end{quote}

\subsection*{You'll want to know: \kn{ವಿವಾಹ}}
\textbf{\kn{ವಿವಾಹ}} — \emph{vivāha} — \textquotedblleft{}a marriage\textquotedblright{}.

\begin{quote}
\kn{ಎಲ್ಲವೂ} \kn{ಸರಿಯಾಗಿ} \kn{ನಡೆದರೆ} \kn{ಮುಂದಿನ} \kn{ತಿಂಗಳು} \kn{ಅವರ} \kn{ವಿವಾಹ} \kn{ನಡೆಯುತ್ತದೆ}.
\end{quote}

\begin{quote}
\emph{ellavū sariyāgi naḍedare mundina tiṅgaḷu avara vivāha naḍeyuttade.}
\end{quote}

\begin{quote}
If all goes well, their marriage takes place next month.
\end{quote}

\kn{ವಿವಾಹ} is the formal word, on cards and in the news; at home people say \kn{ಮದುವೆ}.

\begin{cousinweb}
\textbf{\kn{ವಿವಾಹ}} (\emph{vivāha}) is borrowed from Sanskrit \emph{vivāha}, \textquotedblleft{}marriage\textquotedblright{}, literally \textquotedblleft{}leading away\textquotedblright{} of the bride (CDIAL 11920).
\end{cousinweb}

\begin{grammarlens}[title={What you've built}]
One more word for telling what happened.
\end{grammarlens}

\subsection*{Your turn}
\begin{itemize}
\raggedright
  \item \emph{Say it:} \emph{vivāha}
  \item \emph{Say it:} \emph{vivāha}, once more
  \item \emph{Say it:} \emph{ellavū sariyāgi naḍedare mundina tiṅgaḷu avara vivāha naḍeyuttade.} — If all goes well, their marriage takes place next month.
  \item \emph{Recall:} say the Kannada for to extinguish, then the Kannada for to rescue, then say \emph{vivāha} again
\end{itemize}

\begin{tcolorbox}[breakable,colback=teal!4,colframe=teal!35!black,title={Before you move on}]
What does \kn{ವಿವಾಹ} mean? (\textquotedblleft{}A marriage\textquotedblright{}.) Say it once more.
\end{tcolorbox}
\section[\texorpdfstring{\kn{ಆರತಕ್ಷತೆ} (āratakṣate)}{ಆರತಕ್ಷತೆ (āratakṣate)}]{\kn{ಆರತಕ್ಷತೆ} (āratakṣate) — a wedding reception}

\label{lesson:KA-C325-arataksate}



\begin{quote}
Before the new one: say the Kannada for to rescue, then the Kannada for a marriage.
\end{quote}

\subsection*{You'll want to know: \kn{ಆರತಕ್ಷತೆ}}
\textbf{\kn{ಆರತಕ್ಷತೆ}} — \emph{āratakṣate} — \textquotedblleft{}a wedding reception\textquotedblright{}.

\begin{quote}
\kn{ಮದುವೆಯ} \kn{ಹಿಂದಿನ} \kn{ಸಂಜೆ} \kn{ಆರತಕ್ಷತೆ} \kn{ಇತ್ತು}.
\end{quote}

\begin{quote}
\emph{maduveya hindina sañje āratakṣate ittu.}
\end{quote}

\begin{quote}
The wedding reception was on the evening before the wedding.
\end{quote}

The evening when guests meet the couple and bless them with rice grains.

\begin{cousinweb}
\textbf{\kn{ಆರತಕ್ಷತೆ}} (\emph{āratakṣate}) is usually traced to \kn{ಆರತಿ}, the waving of a lamp, and \kn{ಅಕ್ಷತೆ}, grains of unbroken rice, from Sanskrit \emph{akṣata}, \textquotedblleft{}unbroken\textquotedblright{} (Monier-Williams).
\end{cousinweb}

\begin{grammarlens}[title={What you've built}]
One more word for telling what happened.
\end{grammarlens}

\subsection*{Your turn}
\begin{itemize}
\raggedright
  \item \emph{Say it:} \emph{āratakṣate}
  \item \emph{Say it:} \emph{āratakṣate}, once more
  \item \emph{Say it:} \emph{maduveya hindina sañje āratakṣate ittu.} — The wedding reception was on the evening before the wedding.
  \item \emph{Recall:} say the Kannada for to rescue, then the Kannada for a marriage, then say \emph{āratakṣate} again
\end{itemize}

\begin{tcolorbox}[breakable,colback=teal!4,colframe=teal!35!black,title={Before you move on}]
What does \kn{ಆರತಕ್ಷತೆ} mean? (\textquotedblleft{}A wedding reception\textquotedblright{}.) Say it once more.
\end{tcolorbox}
\section[\texorpdfstring{\kn{ವರದಕ್ಷಿಣೆ} (varadakṣiṇe)}{ವರದಕ್ಷಿಣೆ (varadakṣiṇe)}]{\kn{ವರದಕ್ಷಿಣೆ} (varadakṣiṇe) — a dowry}

\label{lesson:KA-C325-varadaksine}



\begin{quote}
Before the new one: say the Kannada for a marriage, then the Kannada for a wedding reception.
\end{quote}

\subsection*{You'll want to know: \kn{ವರದಕ್ಷಿಣೆ}}
\textbf{\kn{ವರದಕ್ಷಿಣೆ}} — \emph{varadakṣiṇe} — \textquotedblleft{}a dowry\textquotedblright{}.

\begin{quote}
\kn{ಆ} \kn{ಕುಟುಂಬ} \kn{ವರದಕ್ಷಿಣೆ} \kn{ಕೇಳಿದ್ದರಿಂದ} \kn{ಮದುವೆ} \kn{ನಿಂತುಹೋಯಿತು}.
\end{quote}

\begin{quote}
\emph{ā kuṭumba varadakṣiṇe kēḷiddarinda maduve nintuhōyitu.}
\end{quote}

\begin{quote}
The wedding was called off because that family asked for a dowry.
\end{quote}

The money or gifts demanded from the bride's family; the practice is against the law.

\begin{cousinweb}
\textbf{\kn{ವರದಕ್ಷಿಣೆ}} (\emph{varadakṣiṇe}) is borrowed from Sanskrit, a compound of \emph{vara}, \textquotedblleft{}a suitor, a bridegroom\textquotedblright{} (CDIAL 11309), and \emph{dakṣiṇā}, \textquotedblleft{}a fee, a gift\textquotedblright{} (CDIAL 6120).
\end{cousinweb}

\begin{grammarlens}[title={What you've built}]
One more word for telling what happened.
\end{grammarlens}

\subsection*{Your turn}
\begin{itemize}
\raggedright
  \item \emph{Say it:} \emph{varadakṣiṇe}
  \item \emph{Say it:} \emph{varadakṣiṇe}, once more
  \item \emph{Say it:} \emph{ā kuṭumba varadakṣiṇe kēḷiddarinda maduve nintuhōyitu.} — The wedding was called off because that family asked for a dowry.
  \item \emph{Recall:} say the Kannada for a marriage, then the Kannada for a wedding reception, then say \emph{varadakṣiṇe} again
\end{itemize}

\begin{tcolorbox}[breakable,colback=teal!4,colframe=teal!35!black,title={Before you move on}]
What does \kn{ವರದಕ್ಷಿಣೆ} mean? (\textquotedblleft{}A dowry\textquotedblright{}.) Say it once more.
\end{tcolorbox}
\section[\texorpdfstring{\kn{ಹರಸು} (harasu)}{ಹರಸು (harasu)}]{\kn{ಹರಸು} (harasu) — to bless}

\label{lesson:KA-C325-harasu}



\begin{quote}
Before the new one: say the Kannada for a wedding reception, then the Kannada for a dowry.
\end{quote}

\subsection*{You'll want to know: \kn{ಹರಸು}}
\textbf{\kn{ಹರಸು}} — \emph{harasu} — \textquotedblleft{}to bless\textquotedblright{}.

\begin{quote}
\kn{ಹಿರಿಯರು} \kn{ಹೊಸ} \kn{ಜೋಡಿಯನ್ನು} \kn{ಹರಸಿದರು}.
\end{quote}

\begin{quote}
\emph{hiriyaru hosa jōḍiyannu harasidaru.}
\end{quote}

\begin{quote}
The elders blessed the newly married couple.
\end{quote}

Past (I …): \kn{ಹರಸಿದೆ}, \emph{haraside}.

The person blessed takes -\kn{ಅನ್ನು}. A blessing or a vow is a \kn{ಹರಕೆ}.

\begin{cousinweb}
\textbf{\kn{ಹರಸು}} (\emph{harasu}) is native Dravidian: DEDR 3951 lists Kannada \emph{parasu}, \textquotedblleft{}to utter a benediction, to bless\textquotedblright{}; old \emph{p-} became \emph{h-}.
\end{cousinweb}

\begin{grammarlens}[title={What you've built}]
One more word for telling what happened.
\end{grammarlens}

\subsection*{Your turn}
\begin{itemize}
\raggedright
  \item \emph{Say it:} \emph{harasu}
  \item \emph{Say it:} \emph{harasu}, once more
  \item \emph{Say it:} \emph{hiriyaru hosa jōḍiyannu harasidaru.} — The elders blessed the newly married couple.
  \item \emph{Recall:} say the Kannada for a wedding reception, then the Kannada for a dowry, then say \emph{harasu} again
\end{itemize}

\begin{tcolorbox}[breakable,colback=teal!4,colframe=teal!35!black,title={Before you move on}]
What does \kn{ಹರಸು} mean? (\textquotedblleft{}To bless\textquotedblright{}.) Say it once more.
\end{tcolorbox}
\section[\texorpdfstring{\kn{ದಿಬ್ಬಣ} (dibbaṇa)}{ದಿಬ್ಬಣ (dibbaṇa)}]{\kn{ದಿಬ್ಬಣ} (dibbaṇa) — a wedding party}

\label{lesson:KA-C325-dibbana}



\begin{quote}
Before the new one: say the Kannada for a dowry, then the Kannada for to bless.
\end{quote}

\subsection*{You'll want to know: \kn{ದಿಬ್ಬಣ}}
\textbf{\kn{ದಿಬ್ಬಣ}} — \emph{dibbaṇa} — \textquotedblleft{}a wedding party\textquotedblright{}.

\begin{quote}
\kn{ಗಂಡಿನ} \kn{ಕಡೆಯ} \kn{ದಿಬ್ಬಣ} \kn{ಸಂಜೆ} \kn{ಊರಿಗೆ} \kn{ಬಂತು}.
\end{quote}

\begin{quote}
\emph{gaṇḍina kaḍeya dibbaṇa sañje ūrige bantu.}
\end{quote}

\begin{quote}
The groom's wedding party reached the village in the evening.
\end{quote}

The procession of the groom's family and friends to the wedding.

\begin{cousinweb}
\textbf{\kn{ದಿಬ್ಬಣ}} (\emph{dibbaṇa}) is native Dravidian: DEDR 3230 records Kannada \emph{dibbaṇa}, \textquotedblleft{}a nuptial procession, the friends and relatives of the bridegroom\textquotedblright{}.
\end{cousinweb}

\begin{grammarlens}[title={What you've built}]
One more word for telling what happened.
\end{grammarlens}

\subsection*{Your turn}
\begin{itemize}
\raggedright
  \item \emph{Say it:} \emph{dibbaṇa}
  \item \emph{Say it:} \emph{dibbaṇa}, once more
  \item \emph{Say it:} \emph{gaṇḍina kaḍeya dibbaṇa sañje ūrige bantu.} — The groom's wedding party reached the village in the evening.
  \item \emph{Recall:} say the Kannada for a dowry, then the Kannada for to bless, then say \emph{dibbaṇa} again
\end{itemize}

\begin{tcolorbox}[breakable,colback=teal!4,colframe=teal!35!black,title={Before you move on}]
What does \kn{ದಿಬ್ಬಣ} mean? (\textquotedblleft{}A wedding party\textquotedblright{}.) Say it once more.
\end{tcolorbox}
